% Bagian: normal-modal-logic
% Bab: tableaux
% Subbagian: countermodels

\documentclass[../../../include/open-logic-section]{subfiles}

\begin{document}

\olfileid[id]{nml}{tab}{cou}
      
\olsection{Model Tandingan dari \usetoken{P}{tableau}}

\begin{explain}
  Pembuktian teorema kelengkapan bukan hanya menunjukkan bahwa jika
  $\Entails !A$, maka $\Proves !A$; pembuktian tersebut juga memberi kita
  metode untuk membangun model tandingan bagi~$!A$ jika $\Entails/ !A$.
  Dalam kasus~$\Log{K}$, metode ini merupakan suatu \emph{prosedur
  keputusan}. Andaikan $\Entails/ !A$. Pembuktian
  \olref[cpl]{prop:complete-tableau} memberi suatu metode untuk membangun
  !!{tableau} lengkap. Metode tersebut bahkan selalu berhenti. Aturan
  proposisional bagi~$\Log{K}$ hanya menambahkan !!{formula}s berprefiks
  dengan kompleksitas yang lebih rendah; dengan kata lain, setiap aturan
  proposisional hanya perlu diterapkan sekali pada suatu cabang bagi setiap
  formula bertanda $\sFmla{S}{!A}[\sigma]$. Prefiks baru hanya dihasilkan
  oleh aturan
  \iftag{prvBox}{$\TRule{\False}{\Box}$}{}\iftag{notprvBox,notprvDiamond}{}{
    dan }\iftag{prvDiamond}{$\TRule{\True}{\Diamond}$}{}, dan
  \iftag{notprvBox,notprvDiamond}{aturan itu}{aturan-aturan itu} juga hanya
  perlu diterapkan sekali (serta menghasilkan satu prefiks baru).
  \iftag{prvBox}{$\TRule{\True}{\Box}$}{}\iftag{notprvBox,notprvDiamond}{}{
    dan }\iftag{prvDiamond}{$\TRule{\False}{\Diamond}$}{}
  \iftag{notprvBox,notprvDiamond}{harus}{harus} mungkin diterapkan beberapa
  kali, tetapi hanya sekali per prefiks, dan hanya berhingga banyak prefiks
  baru yang dihasilkan. Jadi, setelah berhingga banyak tahap, konstruksi
  tersebut menghasilkan cabang tertutup atau cabang lengkap.

  Setelah suatu tableau dengan cabang terbuka yang lengkap dibangun,
  pembuktian \olref[cpl]{thm:tableau-completeness} memberi kita suatu model
  eksplisit yang memenuhi himpunan awal !!{formula}s berprefiks. Jadi, bukan
  hanya benar bahwa jika $\Gamma \Entails !A$, maka terdapat !!{tableau}
  tertutup dan $\Gamma \Proves !A$; jika kita mencari !!{tableau} tertutup
  dengan cara yang tepat dan memperoleh !!{tableau} ``lengkap'', kita bukan
  hanya akan mengetahui bahwa $\Gamma \Entails/ !A$, melainkan juga benar-benar
  dapat membangun suatu model tandingan.
\end{explain}

\iftag{prvBox}{
\begin{ex}
  Kita mengetahui bahwa $\Proves/ \Box(p \lor q) \lif (\Box p \lor \Box
  q)$. Konstruksi suatu tableau dimulai dengan:
  \begin{oltableau}
    [\pFmla{\False}{\Box(p \lor q) \lif (\Box p \lor \Box q)}{1},
      just = \TAss, checked
      [\pFmla{\True}{\Box(p \lor q)}{1},
        just = {\TRule{\False}{\lif}[1]}, 
        [\pFmla{\False}{\Box p \lor \Box q}{1},
          just = {\TRule{\False}{\lif}[1]}, checked
          [\pFmla{\False}{\Box p}{1},
            just = {\TRule{\False}{\lor}[3]}, checked
            [\pFmla{\False}{\Box q}{1},
              just = {\TRule{\False}{\lor}[3]}, checked
              [\pFmla{\False}{p}{1.1}, 
                just = {\TRule{\False}{\Box}[4]}, checked
                [\pFmla{\False}{q}{1.2}, 
                  just = {\TRule{\False}{\Box}[5]}, checked
                ]
              ]
            ]
          ]
        ]
      ]
    ]
  \end{oltableau}
  Tentu saja !!{tableau} tersebut belum selesai. Pada langkah berikutnya,
  kita membahas satu-satunya baris tanpa tanda centang, yakni !!{formula}
  berprefiks $\sFmla{\True}{\Box(p \lor q)}[1]$ pada baris~$2$. Konstruksi
  tableau lengkap memerintahkan kita untuk menerapkan aturan
  $\TRule{\True}{\Box}$ bagi setiap prefiks yang sudah digunakan pada cabang,
  yakni bagi $1.1$ dan $1.2$ sekaligus:
  \begin{oltableau}
    [\pFmla{\False}{\Box(p \lor q) \lif (\Box p \lor \Box q)}{1},
      just = \TAss, checked
      [\pFmla{\True}{\Box(p \lor q)}{1},
        just = {\TRule{\False}{\lif}[1]}, 
        [\pFmla{\False}{\Box p \lor \Box q}{1},
          just = {\TRule{\False}{\lif}[1]}, checked
          [\pFmla{\False}{\Box p}{1},
            just = {\TRule{\False}{\lor}[3]}, checked
            [\pFmla{\False}{\Box q}{1},
              just = {\TRule{\False}{\lor}[3]}, checked
              [\pFmla{\False}{p}{1.1}, 
                just = {\TRule{\False}{\Box}[4]}, checked
                [\pFmla{\False}{q}{1.2}, 
                  just = {\TRule{\False}{\Box}[5]}, checked
                  [\pFmla{\True}{p \lor q}{1.1}, 
                    just = {\TRule{\True}{\Box}[2]}
                    [\pFmla{\True}{p \lor q}{1.2}, 
                      just = {\TRule{\True}{\Box}[2]}
                    ]
                  ]
                ]
              ]
            ]
          ]
        ]
      ]
    ]
  \end{oltableau}
  Sekarang baris 2, 8, dan 9 tidak mempunyai tanda centang. Namun, tidak ada
  prefiks baru yang telah ditambahkan, sehingga kita menerapkan
  $\TRule{\True}{\lor}$ pada baris~8 dan~9 pada semua cabang yang dihasilkan
  (selama cabang-cabang itu tidak tertutup):
  \begin{oltableau}
    [\pFmla{\False}{\Box(p \lor q) \lif (\Box p \lor \Box q)}{1},
      just = \TAss, checked
      [\pFmla{\True}{\Box(p \lor q)}{1},
        just = {\TRule{\False}{\lif}[1]}, checked
        [\pFmla{\False}{\Box p \lor \Box q}{1},
          just = {\TRule{\False}{\lif}[1]}, checked
          [\pFmla{\False}{\Box p}{1},
            just = {\TRule{\False}{\lor}[3]}, checked
            [\pFmla{\False}{\Box q}{1},
              just = {\TRule{\False}{\lor}[3]}, checked
              [\pFmla{\False}{p}{1.1}, 
                just = {\TRule{\False}{\Box}[4]}, checked
                [\pFmla{\False}{q}{1.2}, 
                  just = {\TRule{\False}{\Box}[5]}, checked
                  [\pFmla{\True}{p \lor q}{1.1}, 
                    just = {\TRule{\True}{\Box}[2]}, checked
                    [\pFmla{\True}{p \lor q}{1.2}, 
                      just = {\TRule{\True}{\Box}[2]}, checked
                      [\pFmla{\True}{p}{1.1},
                        just = {\TRule{\True}{\lor}[8]}, checked, close
                      ]
                      [\pFmla{\True}{q}{1.1},
                        just = {\TRule{\True}{\lor}[8]}, checked
                        [\pFmla{\True}{p}{1.2},
                          just = {\TRule{\True}{\lor}[9]}, checked]
                        [\pFmla{\True}{q}{1.2},
                          just = {\TRule{\True}{\lor}[9]}, checked, close]
                      ]
                    ]
                  ]
                ]
              ]
            ]
          ]
        ]
      ]
    ]
  \end{oltableau}
  Tersisa satu cabang terbuka, dan cabang itu lengkap. Dari cabang tersebut,
  kita mendefinisikan model dengan dunia $W = \{1, 1.1, 1.2\}$ (satu-satunya
  prefiks yang muncul pada cabang terbuka), relasi aksesibilitas
  $R = \{\tuple{1, 1.1}, \tuple{1, 1.2}\}$, serta penetapan
  $V(p) = \{1.2\}$ (karena baris~11 memuat
  $\sFmla{\True}{p}[1.2]$) dan $V(q) = \{1.1\}$ (karena baris~10 memuat
  $\sFmla{\True}{q}[1.1]$). Model itu digambarkan dalam
  \olref{fig:counter-Box}, dan anda dapat memeriksa bahwa model itu merupakan
  model tandingan bagi $\Box(p \lor q) \lif (\Box p \lor \Box q)$.
  \begin{figure}
  \begin{center}
    \begin{tikzpicture}[modal]
      \node[world] (w1) [label={[align=right]right:\mFalse{p}\\ \mFalse{q}}]{$1$}; 
      \node[world] (w2) [label={[align=right]right:\mFalse{p}\\ \mTrue{q}},
        above left=of w1]{$1.1$}; 
      \node[world] (w3) [label={[align=right]right:\mTrue{p}\\ \mFalse{q}},
        above right=of w1] {$1.2$};
      \draw[->] (w1) to (w2);
      \draw[->] (w1) to (w3);
    \end{tikzpicture}
  \end{center}
  \caption{Model tandingan bagi $\Box(p \lor q) \lif (\Box p \lor \Box
    q)$.}
  \ollabel{fig:counter-Box}
\end{figure}
\end{ex}          
}
{
\begin{ex}
  Kita mengetahui bahwa $\Proves/ (\Diamond p \land \Diamond q) \lif
  \Diamond(p \land q)$. Konstruksi suatu tableau dimulai dengan:
  \begin{oltableau}
    [\pFmla{\False}{(\Diamond p \land \Diamond q) \lif \Diamond(p \land q)}{1},
      just = \TAss, checked
      [\pFmla{\True}{\Diamond p \land \Diamond q}{1},
        just = {\TRule{\False}{\lif}[1]}, checked
        [\pFmla{\False}{\Diamond(p \land q)}{1},
          just = {\TRule{\False}{\lif}[1]}
          [\pFmla{\True}{\Diamond p}{1},
            just = {\TRule{\True}{\land}[2]}, checked
            [\pFmla{\True}{\Diamond q}{1},
              just = {\TRule{\True}{\land}[2]}, checked
              [\pFmla{\True}{p}{1.1}, 
                just = {\TRule{\True}{\Diamond}[4]}, checked
                [\pFmla{\True}{q}{1.2}, 
                  just = {\TRule{\True}{\Diamond}[5]}, checked
                ]
              ]
            ]
          ]
        ]
      ]
    ]
  \end{oltableau}
  Tentu saja !!{tableau} tersebut belum selesai. Pada langkah berikutnya,
  kita membahas satu-satunya baris tanpa tanda centang, yakni !!{formula}
  berprefiks $\sFmla{\False}{\Diamond(p \land q)}[1]$ pada baris~$3$.
  Konstruksi tableau lengkap memerintahkan kita untuk menerapkan aturan
  $\TRule{\False}{\Diamond}$ bagi setiap prefiks yang sudah digunakan pada
  cabang, yakni bagi $1.1$ dan $1.2$ sekaligus:
  \begin{oltableau}
    [\pFmla{\False}{(\Diamond p \land \Diamond q) \lif \Diamond(p \land q)}{1},
      just = \TAss, checked
      [\pFmla{\True}{\Diamond p \land \Diamond q}{1},
        just = {\TRule{\False}{\lif}[1]}, checked
        [\pFmla{\False}{\Diamond (p \land q)}{1},
          just = {\TRule{\False}{\lif}[1]}
          [\pFmla{\True}{\Diamond p}{1},
            just = {\TRule{\True}{\land}[2]}, checked
            [\pFmla{\True}{\Diamond q}{1},
              just = {\TRule{\True}{\land}[2]}, checked
              [\pFmla{\True}{p}{1.1}, 
                just = {\TRule{\True}{\Diamond}[4]}, checked
                [\pFmla{\True}{q}{1.2}, 
                  just = {\TRule{\True}{\Diamond}[5]}, checked
                  [\pFmla{\False}{p \land q}{1.1}, 
                    just = {\TRule{\False}{\Diamond}[3]}
                    [\pFmla{\False}{p \land q}{1.2}, 
                      just = {\TRule{\False}{\Diamond}[3]}
                    ]
                  ]
                ]
              ]
            ]
          ]
        ]
      ]
    ]
  \end{oltableau}
  Sekarang baris 3, 8, dan 9 tidak mempunyai tanda centang. Namun, tidak ada
  prefiks baru yang telah ditambahkan, sehingga kita menerapkan
  $\TRule{\False}{\land}$ pada baris~8 dan~9 pada semua cabang yang dihasilkan
  (selama cabang-cabang itu tidak tertutup):
  \begin{oltableau}
    [\pFmla{\False}{(\Diamond p \land \Diamond q) \lif \Diamond(p \land q)}{1},
      just = \TAss, checked
      [\pFmla{\True}{\Diamond p \land \Diamond q}{1},
        just = {\TRule{\False}{\lif}[1]}, checked
        [\pFmla{\False}{\Diamond(p \land q)}{1},
          just = {\TRule{\False}{\lif}[1]}
          [\pFmla{\True}{\Diamond p}{1},
            just = {\TRule{\True}{\land}[2]}, checked
            [\pFmla{\True}{\Diamond q}{1},
              just = {\TRule{\True}{\land}[2]}, checked
              [\pFmla{\True}{p}{1.1}, 
                just = {\TRule{\True}{\Diamond}[4]}, checked
                [\pFmla{\True}{q}{1.2}, 
                  just = {\TRule{\True}{\Diamond}[5]}, checked
                  [\pFmla{\False}{p \land q}{1.1}, 
                    just = {\TRule{\False}{\Diamond}[3]}, checked
                    [\pFmla{\False}{p \land q}{1.2}, 
                      just = {\TRule{\False}{\Diamond}[3]}, checked
                      [\pFmla{\False}{p}{1.1},
                        just = {\TRule{\False}{\land}[8]}, checked, close
                      ]
                      [\pFmla{\False}{q}{1.1},
                        just = {\TRule{\False}{\land}[8]}, checked
                        [\pFmla{\False}{p}{1.2},
                          just = {\TRule{\False}{\land}[9]}, checked]
                        [\pFmla{\False}{q}{1.2},
                          just = {\TRule{\False}{\land}[9]}, checked, close]
                      ]
                    ]
                  ]
                ]
              ]
            ]
          ]
        ]
      ]
    ]
  \end{oltableau}
  Tersisa satu cabang terbuka, dan cabang itu lengkap. Dari cabang tersebut,
  kita mendefinisikan model dengan dunia $W = \{1, 1.1, 1.2\}$ (satu-satunya
  prefiks yang muncul pada cabang terbuka), relasi aksesibilitas
  $R = \{\tuple{1, 1.1}, \tuple{1, 1.2}\}$, serta penetapan
  $V(p) = \{1.1\}$ (karena baris~6 memuat
  $\sFmla{\True}{p}[1.1]$) dan $V(q) = \{1.2\}$ (karena baris~7 memuat
  $\sFmla{\True}{q}[1.2]$). Model itu digambarkan dalam
  \olref{fig:counter-Diamond}, dan anda dapat memeriksa bahwa model itu
  merupakan model tandingan bagi
  $(\Diamond p \land \Diamond q) \lif \Diamond (p \land q)$.
  \begin{figure}
  \begin{center}
    \begin{tikzpicture}[modal]
      \node[world] (w1) [label={[align=right]right:\mFalse{p}\\ \mFalse{q}}]{$1$}; 
      \node[world] (w2) [label={[align=right]right:\mTrue{p}\\ \mFalse{q}},
        above left=of w1]{$1.1$}; 
      \node[world] (w3) [label={[align=right]right:\mFalse{p}\\ \mTrue{q}},
        above right=of w1] {$1.2$};
      \draw[->] (w1) to (w2);
      \draw[->] (w1) to (w3);
    \end{tikzpicture}
  \end{center}
  \caption{Model tandingan bagi $(\Diamond p \land \Diamond q) \lif
    \Diamond (p \land q)$.}
  \ollabel{fig:counter-Diamond}
\end{figure}
\end{ex}          
}

\end{document}
